%------------------------------------------------------------------------------
% 4.9 ملخص الفصل
% المالك: TM (ضمن مهمة ٢.٢.٧)
% الغرض: تكثيف القرارات التصميمية (لا سرد العناوين) وربطها بمخرجات الفصل
%   الثالث وتمهيد الفصل الخامس الذي ينفذ النظام على هذه الأسس.
%
% عناصر مطلوبة:
%   [x] قائمة النتائج التصميمية: 7 نقاط جوهرية
%   [x] جملة الانتقال إلى الفصل الخامس
%
% حالة الإنجاز: أُنجز (2026-11-04) — أحمد المتوكل؛ وأُعيدت المراجعة
%   الختامية (2026-11-12) بعد اكتمال الأقسام 4.1–4.8: حُدّثت نقطة
%   التخزين لمخطط القسم 4.4 المراجع (تسعة وعشرون جدولاً)، ونقطة الأمان
%   لصياغة القسم 4.5 النهائية، ونقطة الواجهات البرمجية لصياغة القسم
%   4.6 النهائية (اثنتا عشرة نقطة وسبعة أنواع أحداث وصيغة خطأ ثلاثية)،
%   وطُوبقت النقاط السبع كلها مع صياغات ملاكها النهائية.
%------------------------------------------------------------------------------
\section{ملخص الفصل}
\label{sec:design-summary}

حوّل هذا الفصل المتطلبات التسعة والعشرين وحالات الاستخدام السبع التي أقرها التحليل إلى قرارات تصميمية قابلة للتنفيذ والتحقق، على المستويين المتكاملين: المعمارية الكلية التي تحكم المكونات وعلاقاتها، والتصاميم التفصيلية التي تحكم كل طبقة على حدة. وأبرز نتائج الفصل التصميمية يمكن تكثيفها في النقاط الآتية:

\begin{itemize}[label=\textcolor{kauprimary}{$\blacktriangleright$},leftmargin=*]
  \item \textbf{معمارية موثقة البدائل:} بنية ثلاثية الطبقات بوكلاء لغويين منسقين برسم بياني حالة (\EN{LangGraph}) ومشرف مركزي، مثبتة بمخطط مكونات \EN{UML} بتسعة مكونات وواجهات موفرة ومطلوبة، وثمانية قرارات معمارية لكل منها بديل مرفوض ومبرر مرتبط بمتطلب محدد (القسم \ref{sec:architecture}).
  \item \textbf{عقود تنفيذية للوكلاء:} خمس عقد بعقود قراءة وكتابة صريحة وسلوك فشل منضبط، وحالة مشتركة من تسعة حقول باختزالي إلحاق، وقواعد توجيه مشروطة تشمل حلقة إصلاح ذاتي بثلاث محاولات، وعشر أدوات مقيدة الأمان مطابقة للنموذج التحليلي في القسم \ref{sec:agent-modeling} (القسم \ref{sec:agents-design}).
  \item \textbf{طبقة عروض ذكية بتصنيف أمني رباعي:} عروض علائقية موصوفة بالعربية تختزل أكثر من 340 جدولاً خاماً إلى ثلاثة عشر عرضاً، بكتالوج يقدم للنموذج وتصنيف أمني يحكم ظهورها لكل دور ودورة حياة موثقة (القسم \ref{sec:smart-views-design}).
  \item \textbf{تخزين مزدوج الفصل:} قراءة فقط من قاعدة بيانات المؤسسة عبر مستخدم معزول وطبقة العروض، وقاعدة منصة (\EN{PostgreSQL 16} مع \EN{pgvector}) بمخطط منطقي مطبَّع من تسعة وعشرين جدولاً للهوية والصلاحيات والعروض والفهرس الدلالي والمستندات والتقارير والتدقيق، ومخزن متجهي مستقل للمعرفة اللوائحية (القسم \ref{sec:db-design}).
  \item \textbf{إطار أمان مبني على الأدوار:} مصفوفة صلاحيات لأربعة أدوار مقابل حالات الاستخدام السبع بمنعٍ صريحٍ مقدم على السماح، وجلسات قابلة للسحب برموز قصيرة الصلاحية، وسجل تدقيق منقَّح يوثق قرارات الوصول وتتبع الوكلاء (القسم \ref{sec:security-design}).
  \item \textbf{كتالوج واجهات برمجية موحد:} اثنتا عشرة نقطة نهاية ببادئة إصدارية واحدة موثقة توثيقاً آلياً بمواصفة \EN{OpenAPI}، وبث خادم للأحداث \EN{(SSE)} بسبعة أنواع معيارية يحوّل رسائل الانتظار إلى تغذية راجعة حية قابلة للاستئناف بعد الانقطاع، وصيغة خطأ موحدة برمز داخلي ثابت ورسالة عربية للمستخدم وتفاصيل تشخيصية تثبت في التدقيق (القسم \ref{sec:api-design}).
  \item \textbf{واجهة عربية أولاً ودليل اكتمال التصميم:} أنماط حوارية مصممة لاتجاه القراءة العربية مع لوحة تحكم إدارية (القسم \ref{sec:ui-design})، ومصفوفة تتبع تربط كل عنصر تصميم بمتطلبه الأصلي دليلَ اكتمالٍ للتغطية (القسم \ref{sec:design-traceability}).
\end{itemize}

وبذلك اكتمل عقد التصميم: كل متطلب وظيفي له موضع تصميم يحققه، وكل قرار معماري له مبرر موثق ومتطلب يخدمه، وكل مكون له عقد وواجهة ومالك تنفيذ. وعلى هذه الأسس ينتقل الفصل الخامس إلى بيئة التنفيذ الفعلية: حاويات \EN{Docker Compose} بست خدمات، وبيئة اختبار \EN{Odoo} جامعية ببيانات واقعية، ونموذج لغوي محلي يعمل عبر \EN{Ollama}، ليبني النظام وفق هذا التصميم ويقيسه على مجموعة تقييم مرجعية بأهداف دقة وزمن استجابة محددة سلفاً.
